\documentclass[a4paper,11pt]{article}
\usepackage[T2A]{fontenc}
\usepackage[utf8]{inputenc}
\usepackage[russian]{babel}
\usepackage{amsmath,amssymb,mathtools}
\usepackage{helvet}
\renewcommand{\familydefault}{\sfdefault}
\usepackage{icomma}
\usepackage{geometry}
\geometry{left=23mm,right=22mm,top=23mm,bottom=26mm}
\usepackage{microtype}
\usepackage{tabularx,booktabs,array}
\usepackage{enumitem}
\usepackage{tikz}
\usepackage{hyperref}
\usepackage{fancyhdr}
\usepackage{xcolor}
\usepackage{setspace}
\definecolor{accent}{HTML}{006B72}
\definecolor{darkaccent}{HTML}{24484C}
\definecolor{ink}{HTML}{242A2E}
\definecolor{muted}{HTML}{66767C}
\definecolor{errorred}{HTML}{A93236}
\definecolor{conceptblue}{HTML}{245E99}
\hypersetup{colorlinks=true,linkcolor=accent,pdfauthor={Лёвин Артём Александрович},pdftitle={Ревью домашней работы — Анастасия Павлова, 05.11.26}}
\setlength{\parindent}{0pt}
\setlength{\parskip}{3pt plus 1pt minus 1pt}
\setlength{\headheight}{14pt}
\renewcommand{\arraystretch}{1.28}
\sloppy
\pagestyle{fancy}
\fancyhf{}
\fancyhead[L]{\footnotesize\color{muted}Ревью домашней работы}
\fancyhead[R]{\footnotesize\color{muted}Анастасия Павлова}
\fancyfoot[R]{\footnotesize\color{muted}\thepage}
\fancyfoot[L]{\tiny\color{muted}Лёвин Артём Александрович эксклюзивно для Анастасии Павловой}
\renewcommand{\headrulewidth}{0.25pt}
\renewcommand{\footrulewidth}{0pt}
\newcommand{\parttitle}[1]{\par\addvspace{4pt}{\color{darkaccent}\large\bfseries #1}\par\vspace{1pt}}
\newcommand{\errorline}[1]{\textcolor{errorred}{\textbf{Ошибка:}} #1}
\newcommand{\idea}[1]{\textcolor{conceptblue}{\textbf{Ключевая идея:}} #1}
\newcommand{\resultline}[1]{\textcolor{accent}{\textbf{Итог:}} #1}
\newcommand{\taskhead}[2]{\phantomsection\label{task:#1}{\Huge\color{accent}\bfseries Задание №#1}\par\vspace{3pt}{\large\color{ink}\bfseries #2}\par\vspace{9pt}\hrule\vspace{12pt}}
\begin{document}
\thispagestyle{empty}
\begin{center}
\vspace*{40mm}
{\large\color{muted}МАТЕМАТИКА}\par\vspace{18mm}
{\Huge\bfseries\color{darkaccent}Ревью домашней работы}\par\vspace{10mm}
{\Large Итоговое повторение: 21 задача}\par\vspace{16mm}
{\LARGE Анастасия Павлова}\par\vspace{10mm}
{\large 5 ноября 2026 г.}\par
\vfill
{\normalsize Лёвин Артём Александрович}\par\vspace{2mm}
{\small\color{muted}эксклюзивно для Анастасии Павловой}\par\vspace{10mm}
\begin{tikzpicture}[x=1cm,y=1cm]
\draw[accent!55,line width=.9pt] (-5,0) .. controls (-2.5,.45) and (-1.7,-.35) .. (0,0) .. controls (1.9,.35) and (3,-.45) .. (5,0);
\end{tikzpicture}
\end{center}
\clearpage
\section*{Сводная таблица ошибок}
\vspace{5pt}
В таблицу включены задания, в которых обнаружена ошибка, неполное решение или отсутствует обязательная часть ответа. Номера ведут к подробным разборам.
\vspace{8pt}
\small
\begin{tabularx}{\linewidth}{@{}>{\centering\arraybackslash}p{9mm} p{24mm} X p{26mm} p{31mm}@{}}
\toprule
\textbf{№} & \textbf{Тема} & \textbf{Суть ошибки} & \textbf{Ваш ответ} & \textbf{Правильный ответ}\\
\midrule
\hyperref[task:4]{4} & Степени & Перепутан обратный показатель; не обоснована граница. & $512$ & $V_{\max}=16$\\
\hyperref[task:5]{5} & Функция & Ошибка при вычислении $3^3$; нуль функции не найден. & $a=3$, $b=-3$, $f(3)=6$ & $a=3$, $b=-3$, $x_0=1$, $f(3)=24$\\
\hyperref[task:9]{9} & Встреча & Сложены времена вместо расстояний; второй ответ отсутствует. & $2$ часа & $2$ часа; $120$ км\\
\hyperref[task:15]{15} & Смеси & Масса раствора после выпаривания ошибочно увеличена. & $6$ кг & $6$ кг\\
\hyperref[task:16]{16} & Сушка & Перепутаны масса сухого вещества и масса изюма. & $40$ кг & $50$ кг\\
\hyperref[task:17]{17} & Проценты & Снижение на $25\%$ применено к первоначальной цене. & $\approx2842$ руб.; $5\%$ & $3000$ руб.; снижение $10\%$\\
\hyperref[task:18]{18} & Вклад & Искажён свободный член уравнения; решение не закончено. & Не указан & $12\%$\\
\hyperref[task:21]{21} & Кредит & Выполнена только классификация схем; расчёты пропущены. & Ставка и сумма не указаны & $7\%$; $1\,410\,000$ руб.\\
\bottomrule
\end{tabularx}
\normalsize
\clearpage
\thispagestyle{empty}
\vspace*{.37\textheight}
\begin{center}
{\Huge\bfseries\color{darkaccent}Разбор ошибок}
\end{center}
\clearpage

\taskhead{4}{Прикладная степенная модель}
\parttitle{Текст задания}
Положительные величины $p$ и $V$ связаны законом $pV^{3/2}=16\,000$. Известно, что $p\ge250$. Найдите наибольшее допустимое $V$ и обоснуйте выбор границы $p$.
\parttitle{Ваше решение и ответ}
В записи верно взято $p=250$, затем получено $V^{3/2}=64$. После возведения в степень $2/3$ записано $V=\sqrt{64^3}=512$. Ответ: $512$.
\parttitle{Локализация ошибки}
\errorline{Неверно преобразована дробная степень при выделении $V$.}

\textbf{Последний правильный шаг:}
\[
250V^{3/2}=16\,000\quad\Longrightarrow\quad V^{3/2}=64.
\]
\textbf{Первый неправильный шаг:} переход от $V^{3/2}=64$ к $V=\sqrt{64^3}$.

\textbf{Почему это неверно.} Чтобы убрать показатель $3/2$, нужно возвести обе части уравнения в обратную степень $2/3$. При этом знаменатель дробного показателя $3$ задаёт степень корня, а числитель $2$ --- степень после извлечения корня. Формула
\[
64^{2/3}=\sqrt[3]{64^2}=(\sqrt[3]{64})^2
\]
не совпадает с $\sqrt{64^3}=64^{3/2}$: во втором выражении показатель оставлен прежним, а не заменён на обратный.

\idea{При равенстве $V^q=A$, где $V>0$ и $q\ne0$, получаем $V=A^{1/q}$. Для максимума $V$ сначала исследуйте, как он зависит от $p$.}
\parttitle{Исправление с последнего правильного шага}
\[
V^{3/2}=64\quad\Longrightarrow\quad
V=64^{2/3}=(\sqrt[3]{64})^2=4^2=16.
\]
\parttitle{Полное правильное решение}
Из исходного закона и $p>0$:
\[
V^{3/2}=\frac{16\,000}{p}.
\]
Когда $p$ увеличивается, положительная дробь $16\,000/p$ уменьшается. При $V>0$ функция $V^{3/2}$ возрастает, значит, само $V$ также уменьшается. Поэтому максимальное $V$ достигается при минимально допустимом $p=250$:
\[
V^{3/2}=\frac{16\,000}{250}=64,
\qquad V=64^{2/3}=16.
\]
\parttitle{Проверка}
При $p=250$ и $V=16$ получаем $250\cdot16^{3/2}=250\cdot64=16\,000$. Оба ограничения выполнены.

\resultline{$V_{\max}=16$.}
\parttitle{Задача на закрепление}
Положительные $p$ и $V$ удовлетворяют $pV^{3/2}=13\,500$, причём $p\ge500$. Найдите наибольшее $V$ и обоснуйте выбор $p$.
\clearpage

\taskhead{5}{График показательной функции}
\parttitle{Текст задания}
Функция $f(x)=a^x+b$ при $a>0$, $a\ne1$ проходит через $(0;-2)$ и $(2;6)$. Найдите $a$, $b$, нуль функции и $f(3)$.
\parttitle{Ваше решение и ответ}
Из точки с абсциссой нуль верно получено $b=-3$. Из второй точки верно найдено $a=3$. Затем записано $f(0)=-3$ и $f(3)=3^3-3=9-3=6$. Нуль функции отдельно не найден.
\parttitle{Локализация ошибки}
\errorline{Неправильно вычислены степени с показателями $0$ и $3$; пропущено обязательное требование найти нуль.}

\textbf{Последний правильный шаг:} $a=3$ и $b=-3$.

\textbf{Первый неправильный шаг:} $f(0)=3^0-3=-3$. Любое положительное основание в нулевой степени равно единице: $3^0=1$, поэтому $f(0)=-2$.

\textbf{Следующая ошибка:} $3^3$ заменено на $9$. Но $3^3=3\cdot3\cdot3=27$, а $9=3^2$.

\idea{Нуль функции --- значение аргумента $x$, при котором $f(x)=0$. Подстановка $x=0$ даёт значение функции на оси $Oy$, но не обязательно её нуль.}
\parttitle{Исправление и полное правильное решение}
Точка $(0;-2)$ принадлежит графику, поэтому
\[
a^0+b=-2\quad\Longrightarrow\quad1+b=-2\quad\Longrightarrow\quad b=-3.
\]
Точка $(2;6)$ также принадлежит графику:
\[
a^2-3=6\quad\Longrightarrow\quad a^2=9.
\]
Из условия $a>0$ следует $a=3$. Значит,
\[
f(x)=3^x-3.
\]
Для поиска нуля приравниваем функцию к нулю:
\[
3^x-3=0\quad\Longrightarrow\quad3^x=3^1\quad\Longrightarrow\quad x=1.
\]
Искомое значение:
\[
f(3)=3^3-3=27-3=24.
\]
\parttitle{Проверка}
$f(0)=1-3=-2$, $f(2)=9-3=6$ --- обе заданные точки подходят. Также $f(1)=3-3=0$.

\resultline{$a=3$, $b=-3$, нуль $x=1$, $f(3)=24$.}
\parttitle{Задача на закрепление}
Функция $g(x)=a^x+b$ при $a>0$, $a\ne1$ проходит через $(0;-3)$ и $(2;12)$. Найдите $a$, $b$, нуль функции и $g(3)$.
\clearpage

\taskhead{9}{Встречное движение}
\parttitle{Текст задания}
Из двух городов на расстоянии $270$ км одновременно навстречу выехали автомобили со скоростями $60$ и $75$ км/ч. Через сколько часов они встретятся и сколько километров проедет первый?
\parttitle{Ваше решение и ответ}
Составлена таблица $v,S,t$ с расстояниями $x/60$ и $x/75$, временами $x$ и $x$. Далее записано $x/60+x/75=270$, затем $135x=270$, откуда $x=2$. Ответ: $2$ часа. Расстояние первого автомобиля не указано.
\parttitle{Локализация ошибки}
\errorline{Перепутаны расстояние и время; неверно составлено уравнение и пропущена часть ответа.}

\textbf{Последняя верная идея:} у автомобилей одинаковое время движения до встречи.

\textbf{Первый неправильный шаг:}
\[
\frac{x}{60}+\frac{x}{75}=270.
\]
Если $x$ означает время в часах, расстояния равны $60x$ и $75x$. Если $x$ означает расстояние, то времена $x/60$ и $x/75$ измеряются в часах и их сумма не может равняться $270$ км. Таким образом, в записи нарушено соответствие единиц измерения.

Кроме того, переход к $135x=270$ из указанного уравнения алгебраически неверен. Число $2$ оказалось правильным по совпадению и нуждается в корректном обосновании.

\idea{При встречном движении без задержек пути автомобилей складываются в исходное расстояние. Их суммарная скорость сближения равна сумме скоростей.}
\parttitle{Исправление и полное правильное решение}
Пусть $t>0$ --- число часов до встречи. Первый автомобиль проедет $60t$ км, второй --- $75t$ км. Вместе они преодолеют $270$ км:
\[
60t+75t=270,\qquad135t=270,\qquad t=2\ \text{ч}.
\]
Путь первого автомобиля:
\[
S_1=60\cdot2=120\ \text{км}.
\]
\parttitle{Проверка}
Второй автомобиль проедет $75\cdot2=150$ км. Сумма $120+150=270$ км совпадает с расстоянием между городами.

\resultline{Встреча через $2$ часа; первый проедет $120$ км.}
\parttitle{Задача на закрепление}
Из городов на расстоянии $360$ км одновременно навстречу выехали автомобили со скоростями $80$ и $100$ км/ч. Найдите время встречи и путь первого автомобиля.
\clearpage

\taskhead{15}{Смешивание и выпаривание}
\parttitle{Текст задания}
Смешали $20$ кг $10\%$-го и $30$ кг $30\%$-го растворов соли. Сколько килограммов воды нужно выпарить, чтобы концентрация соли стала $25\%$?
\parttitle{Ваше решение и ответ}
Верно найдена масса соли: $2+9=11$ кг. Далее записано $11=0{,}25(50+x)$, откуда $x=-6$. В конце взят модуль отрицательного числа и получен ответ $6$ кг.
\parttitle{Локализация ошибки}
\errorline{При выпаривании вода уходит, поэтому общая масса раствора уменьшается.}

\textbf{Последний правильный шаг:} масса соли равна $11$ кг, начальная масса смеси равна $50$ кг.

\textbf{Первый неправильный шаг:}
\[
11=0{,}25(50+x).
\]
Величина $x$ обозначает массу выпаренной воды. Следовательно, масса смеси после выпаривания составляет $50-x$, а не $50+x$. Переход к $|{-6}|=6$ не исправляет неверно составленную математическую модель: отрицательный результат сообщает о несогласованности выбранного знака и смысла $x$.

\idea{При испарении чистой воды масса соли сохраняется, а масса раствора уменьшается. Процентная концентрация равна массе соли, делённой на массу всего раствора.}
\parttitle{Исправление с последнего правильного шага}
Пусть $x$ кг --- масса выпаренной воды. Тогда
\[
\frac{11}{50-x}=0{,}25,\qquad 0\le x<50.
\]
Умножаем на положительную массу оставшегося раствора:
\[
11=0{,}25(50-x)=12{,}5-0{,}25x.
\]
Отсюда
\[
0{,}25x=1{,}5,\qquad x=6.
\]
\parttitle{Полное правильное решение и проверка}
Масса соли до выпаривания: $20\cdot0{,}10+30\cdot0{,}30=2+9=11$ кг. Масса смеси $20+30=50$ кг. После удаления $6$ кг воды остаётся $44$ кг смеси, содержащей прежние $11$ кг соли. Проверяем долю соли:
\[
\frac{11}{44}=\frac14=25\%.
\]
\resultline{Нужно выпарить $6$ кг воды. Числовой ответ совпал, но способ его получения требовал исправления.}
\parttitle{Задача на закрепление}
Смешали $20$ кг $15\%$-го и $40$ кг $30\%$-го растворов соли. Сколько килограммов чистой воды нужно выпарить, чтобы концентрация соли стала $30\%$?
\clearpage

\taskhead{16}{Высушивание}
\parttitle{Текст задания}
Виноград содержит $75\%$ воды. После высушивания её стало $20\%$. Сколько килограммов изюма получится из $160$ кг винограда, если сухое вещество сохраняется?
\parttitle{Ваше решение и ответ}
Верно найдена масса сухого вещества: $160\cdot0{,}25=40$ кг. В таблице для готового изюма записана масса $50$ кг и доля сухого вещества $0{,}8$, но итоговый ответ дан как $40$ кг.
\parttitle{Локализация ошибки}
\errorline{Масса сухого вещества принята за массу готового продукта.}

\textbf{Последний правильный шаг:} $40$ кг --- масса сухого вещества, которая не меняется.

\textbf{Первый неправильный (итоговый) шаг:} объявление $40$ кг массой изюма. В изюме по условию ещё содержится $20\%$ воды. Поэтому сухое вещество составляет только $80\%$ массы изюма, а не все $100\%$.

\idea{При сушке постоянна масса сухого вещества. Изменяется общая масса из-за потери воды.}
\parttitle{Исправление и полное правильное решение}
Исходная доля сухого вещества: $100\%-75\%=25\%$. Значит, его масса:
\[
160\cdot0{,}25=40\ \text{кг}.
\]
Пусть $m$ кг --- масса изюма. В изюме сухого вещества $100\%-20\%=80\%$, то есть $0{,}8m$ кг. Оно сохранилось:
\[
0{,}8m=40\quad\Longrightarrow\quad m=\frac{40}{0{,}8}=50\ \text{кг}.
\]
\parttitle{Проверка}
В $50$ кг изюма содержится $10$ кг воды и $40$ кг сухого вещества. Изначально воды было $160-40=120$ кг; после сушки удалилось $110$ кг воды. Все массы согласованы.

\resultline{Получится $50$ кг изюма.}
\parttitle{Задача на закрепление}
Свежие ягоды массой $180$ кг содержат $80\%$ воды. После сушки содержание воды стало $25\%$. Найдите массу высушенных ягод.
\clearpage

\taskhead{17}{Последовательные проценты}
\parttitle{Текст задания}
Цену сначала повысили на $20\%$, затем снизили полученную цену на $25\%$. Итоговая цена --- $2700$ руб. Найдите начальную цену и общее процентное изменение.
\parttitle{Ваше решение и ответ}
Записано повышение до $120\%$, затем итоговая стоимость ошибочно представлена как $0{,}95x=2700$. Отсюда получено $x\approx2842$ руб. и сделан вывод о снижении на $5\%$.
\parttitle{Локализация ошибки}
\errorline{Второй процент взят от первоначальной цены вместо цены после повышения.}

\textbf{Последний правильный шаг:} после повышения на $20\%$ цена становится $1{,}2x$.

\textbf{Первый неправильный шаг:} замена двух последовательных изменений коэффициентом $0{,}95$. Вторые $25\%$ вычисляются от $1{,}2x$. При снижении остаётся $75\%$ именно этой новой суммы, значит коэффициенты перемножаются: $1{,}2\cdot0{,}75=0{,}9$.

Нельзя складывать или вычитать $20\%$ и $25\%$ как изменения одной и той же базы: база второго изменения уже иная.

\idea{Для последовательных изменений умножайте коэффициенты: рост на $r\%$ даёт множитель $1+r/100$, снижение на $s\%$ --- множитель $1-s/100$.}
\parttitle{Исправление и полное правильное решение}
Обозначим первоначальную цену через $x>0$ руб. После первого изменения имеем $1{,}2x$. После второго изменения:
\[
1{,}2x\cdot(1-0{,}25)=1{,}2x\cdot0{,}75=0{,}9x.
\]
По условию
\[
0{,}9x=2700\quad\Longrightarrow\quad x=\frac{2700}{0{,}9}=3000\ \text{руб}.
\]
Общий множитель равен $0{,}9$, поэтому итоговая цена составляет $90\%$ первоначальной. Снижение равно $100\%-90\%=10\%$.
\parttitle{Проверка}
$3000\cdot1{,}2=3600$ руб.; $3600\cdot0{,}75=2700$ руб. Оба этапа соответствуют условию.

\resultline{Начальная цена $3000$ руб.; суммарное снижение $10\%$.}
\parttitle{Задача на закрепление}
Цену товара повысили на $30\%$, затем новую цену снизили на $20\%$. Итоговая цена --- $3120$ руб. Найдите первоначальную цену и суммарное процентное изменение.
\clearpage

\taskhead{18}{Вклад с пополнением}
\parttitle{Текст задания}
Внесли $500\,000$ руб. под постоянную годовую ставку с капитализацией. В конце первого года после начисления процентов внесли ещё $200\,000$ руб. После второго года стало $851\,200$ руб. Найдите ставку.
\parttitle{Ваше решение и ответ}
Обозначен годовой множитель $k$. Получено верное квадратное уравнение $500k^2+200k-851{,}2=0$ (в тысячах рублей). Затем уравнение умножено на $10$, но в одной из следующих строк свободный член заменён с $-8512$ на $-6512$. По этому неверному числу начато вычисление дискриминанта. Окончательного значения ставки нет.
\parttitle{Локализация ошибки}
\errorline{При переписывании уравнения изменено число $8512$; вычисления не доведены до ставки.}

\textbf{Последний правильный шаг:}
\[
5000k^2+2000k-8512=0.
\]
\textbf{Первый неправильный шаг:} в формуле дискриминанта использован свободный член $-6512$ вместо $-8512$.

Поэтому вычисленное значение дискриминанта относится уже к другому уравнению. Даже правильно извлечённый квадратный корень не дал бы нужную ставку. В записи также требуется переход от множителя $k$ к процентной ставке $r$.

\idea{Пополнение, внесённое после процентов первого года, участвует только в начислении за второй год. Годовой множитель $k=1+r/100$ положителен.}
\parttitle{Исправление и полное правильное решение}
Пусть $r$ --- ставка в процентах, $k=1+r/100$. После первого года на счёте $500\,000k$ руб. Сразу после пополнения станет $500\,000k+200\,000$ руб. За второй год эта сумма умножится на $k$:
\[
(500\,000k+200\,000)k=851\,200.
\]
После деления на $100$ и переноса:
\[
5000k^2+2000k-8512=0.
\]
Дискриминант равен
\[
D=2000^2-4\cdot5000\cdot(-8512)
=4\,000\,000+170\,240\,000=174\,240\,000=13\,200^2.
\]
Корни:
\[
k=\frac{-2000\pm13\,200}{10\,000}
=1{,}12\quad\text{или}\quad-1{,}52.
\]
Отрицательный корень невозможен для годового множителя вклада. Значит, $k=1{,}12$ и
\[
r=(1{,}12-1)\cdot100\%=12\%.
\]
\parttitle{Проверка}
После первого года: $500\,000\cdot1{,}12=560\,000$ руб. После пополнения: $760\,000$ руб. После второго года: $760\,000\cdot1{,}12=851\,200$ руб.

\resultline{Годовая процентная ставка равна $12\%$.}
\parttitle{Задача на закрепление}
На вклад внесли $200\,000$ руб. под постоянную годовую ставку с капитализацией. В конце первого года после начисления процентов добавили $100\,000$ руб. В конце второго года стало $408\,000$ руб. Найдите ставку.
\clearpage

\taskhead{21}{Выбор схемы кредитования и ограничения}
\parttitle{Текст задания}
Назовите схему: (а) ежегодный одинаковый полный платёж; (б) ежегодное равное уменьшение основного долга; (в) ежегодное уменьшение остатка долга на $25\%$. Для (б) рассмотрите кредит $1\,200\,000$ руб. на $4$ года с целой годовой ставкой. Первый платёж $\le390\,000$ руб., последний $\ge320\,000$ руб. Найдите допустимые ставки и сумму выплат.
\parttitle{Ваше решение и ответ}
В записи указано: (а) аннуитетная, (б) дифференцированная, (в) произвольная схема. Расчёт ставки и суммы выплат отсутствует.
\parttitle{Локализация ошибки}
\errorline{Не выполнена вычислительная часть задания.}

\textbf{Последний правильный шаг:} для случая (б) выбран дифференцированный способ погашения, при котором основной долг уменьшается равными частями.

\textbf{Первый неполный шаг:} после определения схемы (б) отсутствует расчёт равной части основного долга и ограничений на первый и последний платежи. Без этих вычислений нельзя найти ставку и полную сумму выплат.

\idea{В дифференцированной схеме одинаковой является доля погашаемого основного долга. Проценты начисляют на остаток перед платежом, поэтому полные платежи уменьшаются.}
\parttitle{Исправление и полное правильное решение}
Классификация: (а) аннуитетная схема; (б) дифференцированная схема; (в) остаток убывает в геометрической прогрессии с множителем $0{,}75$. Само условие (в) не определяет ни аннуитетный, ни дифференцированный способ: требуется дополнительное правило для платежей.

Пусть $r$ --- целая ставка в процентах, $q=r/100$. Ежегодная часть основного долга:
\[
\frac{1\,200\,000}{4}=300\,000\ \text{руб}.
\]
Перед платежами остаток основного долга равен соответственно $1\,200\,000$, $900\,000$, $600\,000$, $300\,000$ руб. Поэтому полные платежи:
\[
\begin{aligned}
P_1&=300\,000+1\,200\,000q,\\
P_2&=300\,000+900\,000q,\\
P_3&=300\,000+600\,000q,\\
P_4&=300\,000+300\,000q.
\end{aligned}
\]
Из ограничения для первого платежа:
\[
300\,000+1\,200\,000q\le390\,000
\quad\Longrightarrow\quad q\le0{,}075.
\]
Из ограничения для последнего:
\[
300\,000+300\,000q\ge320\,000
\quad\Longrightarrow\quad q\ge\frac1{15}\approx0{,}0667.
\]
Значит, $6{,}67\%\lesssim r\le7{,}5\%$. С учётом целочисленности ставки возможно только $r=7\%$, то есть $q=0{,}07$.

Сумма всех выплат состоит из возврата основного долга и процентов:
\[
S=1\,200\,000+(1\,200\,000+900\,000+600\,000+300\,000)\cdot0{,}07
=1\,410\,000\ \text{руб}.
\]
\parttitle{Проверка}
При $7\%$ первый платёж равен $384\,000$ руб. ($\le390\,000$), последний --- $321\,000$ руб. ($\ge320\,000$). Полные платежи: $384\,000$, $363\,000$, $342\,000$, $321\,000$ руб. Их сумма $1\,410\,000$ руб.

\resultline{Единственная целая ставка --- $7\%$; сумма всех выплат --- $1\,410\,000$ руб.}
\parttitle{Задача на закрепление}
По дифференцированной схеме кредит $900\,000$ руб. погашают за $3$ года равными частями основного долга. Годовая ставка --- целое число процентов. Первый платёж не больше $390\,000$ руб., последний не меньше $320\,000$ руб. Найдите все допустимые ставки и сумму всех выплат для каждой ставки.
\end{document}
